\documentclass{article}
\usepackage[utf8]{inputenc}
\usepackage{a4wide}

\begin{document}

\section*{แลกเหรียญ}

กำหนดให้อาร์เรย์จำนวนเต็มบวก\texttt{coins}จำนวน\texttt{N}ค่า

ซึ่งแสดงถึงเหรียญที่มีราคาต่างกัน และจำนวนเต็ม\texttt{amount}แสดงถึงจำนวนเงินทั้งหมด

ให้ตอบจำนวนเหรียญ\ul{\textbf{น้อยที่สุด}}ที่ต้องใช้เพื่อสร้างจำนวนเงิน\texttt{amount}นั้น

หากจำนวนเงินนั้นไม่สามารถสร้างขึ้นจากการรวมเหรียญใดๆ ได้ ให้ตอบ -1\\



ทั้งนี้ให้ถือว่ามีจำนวนเหรียญแต่ละประเภทมากมายไม่จำกัด\\



\textbf{\hfill\break

}



\textbf{ตัวอย่าง}\\



\textbf{\hfill\break

}



{\def\LTcaptype{none} % do not increment counter

\begin{longtable}[]{@{}

  >{\raggedright\arraybackslash}p{(\linewidth - 4\tabcolsep) * \real{0.3333}}

  >{\raggedright\arraybackslash}p{(\linewidth - 4\tabcolsep) * \real{0.3333}}

  >{\raggedright\arraybackslash}p{(\linewidth - 4\tabcolsep) * \real{0.3333}}@{}}

\toprule\noalign{}

\begin{minipage}[b]{\linewidth}\raggedright

ข้อมูลเข้า

\end{minipage} & \begin{minipage}[b]{\linewidth}\raggedright

ข้อมูลออก

\end{minipage} & \begin{minipage}[b]{\linewidth}\raggedright

คำอธิบาย

\end{minipage} \\

\midrule\noalign{}

\endhead

\bottomrule\noalign{}

\endlastfoot

\begin{minipage}[t]{\linewidth}\raggedright

5 29\\

1 4 7 9 10\\

\strut

\end{minipage} & 3 & \begin{minipage}[t]{\linewidth}\raggedright

\textbf{coins = {[}1 4 7 9 10{]}amount = 29\\

}29 = 9 + 10 + 10\\

\strut

\end{minipage} \\

\end{longtable}

}



\subsection*{Input}
\textbf{บรรทัดที่ 1}เป็นจำนวนเต็ม

คือค่า\texttt{N}และค่าจำนวนเต็ม\texttt{amount}เว้นด้วยช่องว่าง



\textbf{บรรทัดที่

2}เป็นชุดเลขจำนวนเต็ม\texttt{N}ค่าสำหรับ\texttt{coins}เว้นด้วยช่องว่าง



\subsection*{Output}
ชุดเลขจำนวนเต็ม 1 ค่าสำหรับคำตอบ คือจำนวนเหรียญน้อยที่สุดที่ต้องใช้\\



\end{document}
